\documentclass[10pt,a4paper]{amsart}
\usepackage[margin=19mm]{geometry}
\usepackage{amsmath,amssymb,mathtools}
\usepackage[scr=rsfs,cal=euler]{mathalfa}
\usepackage{xfrac}
\usepackage[unicode,hidelinks,bookmarks=false]{hyperref}
\newcommand{\ie}{i.e.\ }
\newcommand{\cf}{cf.\ }
\newcommand{\resp}{respectively\ }
\newcommand{\Subsection}{\subsection}
\providecommand{\nobreakdash}{\nobreak\hbox{-}}
\setlength{\emergencystretch}{2em}
\multlinegap0pt
\usepackage{bm}
\usepackage{amssymb}
\usepackage{tikz}
\usepackage{enumerate}
\usepackage[shortlabels]{enumitem}
\usepackage{csquotes}
\usepackage{float}
\usepackage{xcolor}
\usepackage{mathtools}
\newtheorem{lemma}{Lemma}
\newtheorem{prop}{Proposition}
\newcommand{\R}{\mathbb R}
\title{Kusner's conjecture: Exact values and linear bounds}
\author{Hong-Jun Ge, Zixiang Xu, Yang Zhou}
\date{Source: arXiv:2606.03987v1 (CC BY 4.0)}
\begin{document}
\maketitle

\noindent\emph{Source and licence.} Kusner's conjecture: Exact values and linear bounds. Version arXiv:2606.03987v1 (CC BY 4.0), distributed by arXiv under the Creative Commons Attribution 4.0 International licence, http://creativecommons.org/licenses/by/4.0/, as recorded in the arXiv licence metadata for this exact version and shown on the abstract page https://arxiv.org/abs/2606.03987v1. Full licence notice: \texttt{licenses/arxiv-2606.03987-CC-BY-4.0.txt}.\par\smallskip

\noindent\textit{Editorial scope.} Counted unit: the article's prop prop:zero-anchored-exact-bridge-log (source0.tex lines 1877--1925). The article's complete original proof of it is delivered with it (source0.tex lines 1926--2313, one \texttt{proof} environment of 388 lines). No mathematics is written, repaired or re-derived: the statement and the proof are the article's own lines, in the article's own notation, and no retained line is substituted or re-flowed.  Retained uncounted as the source-local context the proof needs: lines 1700--1755.  Nothing was omitted from any retained span, so the package carries no omission record.  No retained line contains a citation, so the wrapper carries no bibliography and no bibliography entry.  No retained line contains a figure environment, an image-inclusion command or drawing code, so no image is delivered and the package declares no asset dependency.  Editorial formatting only, in addition: an AMS \texttt{amsart} preamble that restates the article's own declarations and macros used by the retained lines, namely the environments \texttt{lemma}, \texttt{prop} on separate counters, together with the standard packages those lines need.  The retained environments print in this document as Lemma 1, Proposition 1 in order of appearance, whereas the article prints the same items with the numbering of its own full document.  The complete native source of the article as distributed in the arXiv e-print for this exact version is bundled unchanged as \texttt{sources/202017/source0.tex}.
\begin{lemma}\label{lem:zero-anchored-primitive-log}
For each $x\in[-1,1]$, define
\[
\psi_x^{(0)}(t):=
\begin{cases}
\dfrac{(x-t)_+^{\,n_0}}{n_0!}
-
\displaystyle\sum_{r=0}^{n_0}
\frac{x^r}{r!}\,
\frac{(-t)_+^{\,n_0-r}}{(n_0-r)!},
& x\ge 0,\\[3ex]
\dfrac{(t-x)_+^{\,n_0}}{n_0!}
-
\displaystyle\sum_{r=0}^{n_0}
\frac{(-x)^r}{r!}\,
\frac{t_+^{\,n_0-r}}{(n_0-r)!},
& x\le 0.
\end{cases}
\]
Then the following hold.

\begin{enumerate}
\item[\textnormal{(i)}]
\(
\operatorname{supp}\psi_x^{(0)}
\subseteq [\min\{x,0\},\max\{x,0\}].
\)

\item[\textnormal{(ii)}]
In the sense of distributions,
\[
\bigl(\psi_x^{(0)}\bigr)^{(2k)}
=
\delta_x-\sum_{r=0}^{2k-1}\frac{(-1)^r x^r}{r!}\,\delta_0^{(r)}.
\]

\item[\textnormal{(iii)}]
For every $\alpha>-1$ and every integer $r\in\{0,1,\dots,2k-1\}$,
\[
\int_{\mathbb R}\psi_x^{(0)}(t)\,\operatorname{sgn}(t)^r|t|^\alpha\,dt
=
\frac{\Gamma(\alpha+1)}{\Gamma(\alpha+2k+1)}
\operatorname{sgn}(x)^r|x|^{\alpha+2k}.
\]
In particular,
\(
\int_{\mathbb R}\psi_x^{(0)}(t)\,dt=\frac{|x|^{2k}}{(2k)!},
\)
and
\[
\int_{\mathbb R}\psi_x^{(0)}(t)\,|t|^\beta\,dt
=
\frac{\Gamma(\beta+1)}{\Gamma(p-2k+1)}|x|^{p-2k}.
\]
\end{enumerate}
\end{lemma}

\begin{prop}\label{prop:zero-anchored-exact-bridge-log}
Let $p\in(4k,4k+2)$ for some integer $k\geq 0$, and set
	\(
	\beta:=p-4k\in(0,2).
	\)
	If $p\geq 1$,
then there exist:
\begin{itemize}
\item a real Hilbert space $\mathcal H_\beta$,
\item a continuous map $\Gamma_0:[-1,1]\to \mathcal H_\beta$,
\item a map $\boldsymbol w_0:[-1,1]\to \mathbb R^{4k}$,
\item symmetric positive semidefinite matrices
\(
M_{+,0},M_{-,0}\in\mathbb R^{4k\times 4k},
\)
\item linear operators $T_a:\mathcal H_\beta\to\mathcal H_\beta$ for $0<a\le 1$,
\end{itemize}
such that the following hold.

\medskip
\noindent
\textnormal{(i)}
For every finite set $X=\{x_1,\dots,x_m\}\subseteq [-1,1]$ and every
$\boldsymbol\lambda=(\lambda_1,\dots,\lambda_m)\in \boldsymbol{1}^{\perp}$,
\[
-\frac12\sum_{i,j=1}^m \lambda_i\lambda_j |x_i-x_j|^p
=
c_{p,k}\left\|\sum_{i=1}^m \lambda_i\Gamma_0(x_i)\right\|_{\mathcal H_\beta}^2 +
\left\|\sum_{i=1}^m \lambda_i M_{+,0}^{1/2}\boldsymbol w_0(x_i)\right\|_2^2
-
\left\|\sum_{i=1}^m \lambda_i M_{-,0}^{1/2}\boldsymbol w_0(x_i)\right\|_2^2,
\]
where
\(
c_{p,k}:=\frac{\Gamma(p+1)}{\Gamma(\beta+1)}.
\)

\medskip
\noindent
\textnormal{(ii)}
For all $0<a\le 1$ and all $x\in[-1,1]$,
\(
\Gamma_0(ax)=T_a\Gamma_0(x),
\)
and for \(\xi\in\mathcal H_\beta,\)
\[
\|T_a\xi\|_{\mathcal H_\beta}=a^{p/2}\|\xi\|_{\mathcal H_\beta}.
\]
\end{prop}

\begin{proof}[Proof of Proposition~\ref{prop:zero-anchored-exact-bridge-log}]
If \(k=0\), then \(1\le p<2\). Since \(|x-y|^p\) is conditionally negative definite on \(\R\), there exist a real Hilbert space \(\mathcal H_p\) and a map \(\Gamma_0:\R\to\mathcal H_p\) such that
\[
-\frac12\sum_{i,j}\lambda_i\lambda_j|x_i-x_j|^p
=
\left\|\sum_i\lambda_i\Gamma_0(x_i)\right\|^2.
\]
We take \(\boldsymbol w_0\) to be the map into \(\mathbb R^0\), and \(M_{+,0},M_{-,0}\) to be the corresponding \(0\times0\) matrices. The second statement follows from homogeneity.

   Assume now that \(k\ge1\). We use signed measures only as a compact notation for combining ordinary integrals and point masses: \(f(t)\,dt\) contributes \(\int \varphi(t)f(t)\,dt\), while \(\delta_y\) contributes \(\varphi(y)\). Let
\[
\mathcal M_0:=
\Bigl\{
\mu:\mu\ \text{is a compactly supported signed Borel measure on }\mathbb R,\
\mu(\mathbb R)=0
\Bigr\}.
\]
Define a symmetric bilinear form on $\mathcal M_0$ by
\[
\langle \mu,\nu\rangle_\beta
:=
-\frac12\iint_{\mathbb R^2}|u-v|^\beta\,d\mu(u)\,d\nu(v).
\]
Since $0<\beta<2$, the kernel $(u,v)\rightarrow |u-v|^\beta$ is conditionally negative
definite on $\mathbb R$, and therefore $\langle\cdot,\cdot\rangle_\beta$ is positive
semidefinite on $\mathcal M_0$. Quotienting by the null space and completing, we
obtain a real Hilbert space $\mathcal H_\beta$.

For $x\in[-1,1]$, let $\psi_x^{(0)}$ be as in
Lemma~\ref{lem:zero-anchored-primitive-log}, and define
\[
\nu_x^{(0)}:=\psi_x^{(0)}(t)\,dt-\frac{|x|^{2k}}{(2k)!}\delta_0.
\]
By Lemma~\ref{lem:zero-anchored-primitive-log}\textnormal{(iii)}, $\nu_x^{(0)}\in \mathcal M_0$.
Define
\(
\Gamma_0(x):=[\nu_x^{(0)}]\in \mathcal H_\beta.
\)
Now let
\(
X=\{x_1,\dots,x_m\}\subseteq [-1,1],
\)
\(
\boldsymbol\lambda=(\lambda_1,\dots,\lambda_m)\in \boldsymbol{1}^{\perp},
\)
and define
\[
Q_{p,X}(\boldsymbol\lambda):=
-\frac12\sum_{i,j=1}^m \lambda_i\lambda_j|x_i-x_j|^p.
\]
Set
\(
F_\lambda^{(0)}(t):=\sum_{i=1}^m \lambda_i\psi_{x_i}^{(0)}(t),\)
\(
\mu_\lambda:=\sum_{i=1}^m \lambda_i\delta_{x_i},
\)
\(
M_r(\boldsymbol\lambda):=\sum_{i=1}^m \lambda_i x_i^r
\)
for \(1\le r\le 2k\),
\[
S_r(\boldsymbol\lambda):=\sum_{i=1}^m \lambda_i\,\operatorname{sgn}(x_i)^r|x_i|^{p-r}
\]
for \(1\le r\le 2k.\)
By Lemma~\ref{lem:zero-anchored-primitive-log}\textnormal{(ii)},
\[
(F_\lambda^{(0)})^{(2k)}
=
\mu_\lambda-\sum_{r=0}^{2k-1}\frac{(-1)^r M_r(\boldsymbol\lambda)}{r!}\delta_0^{(r)}.
\]
Since $\boldsymbol\lambda\in\boldsymbol{1}^{\perp}$, we have $M_0(\boldsymbol\lambda)=0$, and therefore
\begin{equation}\label{eq:primitive-distribution-log}
(F_\lambda^{(0)})^{(2k)}
=
\mu_\lambda-\sum_{r=1}^{2k-1}\frac{(-1)^r M_r(\boldsymbol\lambda)}{r!}\delta_0^{(r)}.
\end{equation}
Define also
\[
\nu_\lambda^{(0)}:=F_\lambda^{(0)}(t)\,dt-\frac{M_{2k}(\boldsymbol\lambda)}{(2k)!}\delta_0.
\]
By Lemma~\ref{lem:zero-anchored-primitive-log}\textnormal{(iii)}, $\nu_\lambda^{(0)}\in\mathcal M_0$, and
\(
\sum_{i=1}^m \lambda_i\Gamma_0(x_i)=[\nu_\lambda^{(0)}].
\)
Let \(\|\cdot\|_{\mathcal H_\beta}\) denote the norm in the Hilbert
space \(\mathcal H_\beta\), induced by the inner product
\[
\langle \mu,\nu\rangle_\beta
:=
-\frac12\iint_{\mathbb R^2}|u-v|^\beta\,d\mu(u)\,d\nu(v).
\]
Since
\(
\sum_{i=1}^m\lambda_i\Gamma_0(x_i)=[\nu_\lambda^{(0)}]\in\mathcal H_\beta,
\)
we have, by the definition of this Hilbert-space norm,
\begin{equation}\label{eq:nu-energy-log}
\left\|\sum_{i=1}^m \lambda_i\Gamma_0(x_i)\right\|_{\mathcal H_\beta}^2
=
-\frac12\iint |u-v|^\beta\,d\nu_\lambda^{(0)}(u)\,d\nu_\lambda^{(0)}(v).
\end{equation}

We now compute $Q_{p,X}(\boldsymbol\lambda)$.

\paragraph{Bulk--bulk term.}
Set
\[
Q^{\mathrm{bulk}}
:=
-\frac12\left\langle
(F_\lambda^{(0)})^{(2k)}\otimes (F_\lambda^{(0)})^{(2k)},
\,|x-y|^p
\right\rangle,
\]
where the tensor product is understood in the usual distributional sense, more precisely,
for distributions \(S,T\) on \(\mathbb R\), the distribution \(S\otimes T\) acts on a
function \(\Phi(x,y)\) by first applying \(T\) in the \(y\)-variable and then \(S\)
in the \(x\)-variable.
Since
\[
\partial_x^{2k}\partial_y^{2k}\left(-\frac12|x-y|^p\right)
=
c_{p,k}\left(-\frac12|x-y|^\beta\right),
\]
integration by parts yields
\[
Q^{\mathrm{bulk}}
=
-\frac{c_{p,k}}2
\iint_{\mathbb R^2}
F_\lambda^{(0)}(u)F_\lambda^{(0)}(v)|u-v|^\beta\,du\,dv.
\]
On the other hand, Lemma~\ref{lem:zero-anchored-primitive-log}\textnormal{(iii)} gives
\[
\int_{\mathbb R}F_\lambda^{(0)}(t)\,dt
=
\frac{M_{2k}(\boldsymbol\lambda)}{(2k)!}\]
and
\[
\int_{\mathbb R}|t|^\beta F_\lambda^{(0)}(t)\,dt
=
\frac{\Gamma(\beta+1)}{\Gamma(p-2k+1)}\,S_{2k}(\boldsymbol\lambda).
\]
Expanding \eqref{eq:nu-energy-log}, and using
\(c_{p,k}:=\frac{\Gamma(p+1)}{\Gamma(\beta+1)},\) we obtain
\begin{equation}\label{eq:bulk-identity-log}
Q^{\mathrm{bulk}}
=
c_{p,k}
\left\|\sum_{i=1}^m \lambda_i \Gamma_0(x_i)\right\|_{\mathcal H_\beta}^2
-
\binom{p}{2k}M_{2k}(\boldsymbol\lambda)S_{2k}(\boldsymbol\lambda).
\end{equation}

\paragraph{Bulk--boundary term.}
For $1\le r\le 2k-1$, define
\(
h_r(x):=\operatorname{sgn}(x)^r |x|^{p-r}.
\) From now on, we use the notation \[ (p)_r:=p(p-1)\cdots(p-r+1) \] for the falling factorial.
Moreover, we use \(\langle \delta_0^{(r)},\cdot\rangle_y\) to denote that the distribution \(\delta_0^{(r)}\) acts on the \(y\)-variable, with \(x\) fixed. Therefore, for \(x\neq0\),
\[
\left\langle \delta_0^{(r)},\,|x-y|^p\right\rangle_y = (-1)^r \partial_y^r |x-y|^p\big|_{y=0}.
\]
Since 
\[
\partial_y^r |x-y|^p\big|_{y=0} = (-1)^r (p)_r \operatorname{sgn}(x)^r |x|^{p-r},
\]
we obtain \[ \left\langle \delta_0^{(r)},\,|x-y|^p\right\rangle_y = (p)_r h_r(x). \]
Using \eqref{eq:primitive-distribution-log}, the total bulk--boundary contribution is
\[
Q^{\mathrm{cross}}
=
-\sum_{r=1}^{2k-1}\frac{(-1)^r M_r(\boldsymbol\lambda)}{r!}\,
\left\langle (F_\lambda^{(0)})^{(2k)},\,(p)_r h_r\right\rangle.
\]
Here the factor $2$ coming from the two cross terms is cancelled by the
prefactor $-\frac12$ in $Q_{p,X}(\boldsymbol\lambda)$.

Notice that $p-r>2k$, so $h_r\in C^{2k}(\mathbb R)$, and
\[
h_r^{(2k)}(x)
=
(p-r)_{2k}\,\operatorname{sgn}(x)^r |x|^{p-r-2k}.
\]
Hence, by integration by parts and
Lemma~\ref{lem:zero-anchored-primitive-log}\textnormal{(iii)} with
$\alpha=p-r-2k>-1$,
\[
\left\langle (F_\lambda^{(0)})^{(2k)},\,h_r\right\rangle
=
S_r(\boldsymbol\lambda).
\]
Since $\frac{(p)_r}{r!}=\binom{p}{r}$, summing over $r=1,\dots,2k-1$ gives
\begin{equation}\label{eq:cross-identity-log}
Q^{\mathrm{cross}}
=
-\sum_{r=1}^{2k-1}(-1)^r\binom{p}{r}M_r(\boldsymbol\lambda)S_r(\boldsymbol\lambda).
\end{equation}

\paragraph{Boundary--boundary term.}
Every boundary--boundary term is of the form
\[
\left\langle \delta_0^{(r)}\otimes \delta_0^{(s)},\,|x-y|^p\right\rangle
=
(-1)^{r+s}\partial_x^r\partial_y^s |x-y|^p\big|_{(0,0)}.
\]
Since $r+s$ is an integer and
\(
r+s\le 4k-2<p,
\)
the derivative of order $r+s$ of $|t|^p$ exists at $0$ and equals $0$. Therefore all
boundary--boundary terms vanish.

Combining \eqref{eq:bulk-identity-log} and \eqref{eq:cross-identity-log}, we conclude that
\begin{equation}\label{eq:exact-bridge-finite-part-log}
Q_{p,X}(\boldsymbol\lambda)
=
c_{p,k}
\left\|\sum_{i=1}^m \lambda_i \Gamma_0(x_i)\right\|_{\mathcal H_\beta}^2
-
\sum_{r=1}^{2k}(-1)^r\binom{p}{r}M_r(\boldsymbol\lambda)S_r(\boldsymbol\lambda).
\end{equation}

We now encode the finite-dimensional correction. Define
\(
u_r(x):=x^r
\)
for \(1\le r\le 2k\),
\(
v_r(x):=\operatorname{sgn}(x)^r |x|^{p-r}\)
for \(1\le r\le 2k.\)
Set
\[
\boldsymbol{w}_0(x):=
\bigl(
u_1(x),\dots,u_{2k}(x),v_1(x),\dots,v_{2k}(x)
\bigr)\in\mathbb R^{4k}.
\]
Define a symmetric matrix $A_{p,k}\in\mathbb R^{4k\times 4k}$ by
\[
A_{p,k}(r,2k+r)=A_{p,k}(2k+r,r)
=
-\frac12(-1)^r\binom{p}{r}
\]
for \(1\le r\le 2k\),
and all other entries equal to $0$. Notice that if
\(
\mathbf m_{w_0}(\boldsymbol\lambda):=\sum_{i=1}^m \lambda_i \boldsymbol w_0(x_i),
\)
then by construction
\[
\mathbf m_{w_0}(\boldsymbol\lambda)^\top A_{p,k}\,\mathbf m_{w_0}(\boldsymbol\lambda)
=
-
\sum_{r=1}^{2k}(-1)^r\binom{p}{r}M_r(\boldsymbol\lambda)S_r(\boldsymbol\lambda).
\]
Thus \eqref{eq:exact-bridge-finite-part-log} becomes
\[
Q_{p,X}(\boldsymbol\lambda)
=
c_{p,k}
\left\|\sum_{i=1}^m \lambda_i \Gamma_0(x_i)\right\|_{\mathcal H_\beta}^2
+
\mathbf m_{w_0}(\boldsymbol\lambda)^\top A_{p,k}\,\mathbf m_{w_0}(\boldsymbol\lambda).
\]
Take a spectral decomposition
\(
A_{p,k}=M_{+,0}-M_{-,0}
\)
with \(
M_{+,0},M_{-,0}\succeq 0.
\)
Then
\[
Q_{p,X}(\boldsymbol\lambda)
=
c_{p,k}
\left\|\sum_{i=1}^m \lambda_i \Gamma_0(x_i)\right\|_{\mathcal H_\beta}^2 +
\left\|\sum_{i=1}^m \lambda_i M_{+,0}^{1/2}\boldsymbol w_0(x_i)\right\|_2^2
-
\left\|\sum_{i=1}^m \lambda_i M_{-,0}^{1/2}\boldsymbol w_0(x_i)\right\|_2^2.
\]
This proves \textnormal{(i)}.


We now prove \textnormal{(ii)}.  We first define the scaling operation at the
level of measures.  For \(0<a\le1\) and \(\mu\in\mathcal M_0\), let
\(\widetilde T_a\mu\) be the signed measure determined by
\[
\int_{\mathbb R}\varphi(t)\,d(\widetilde T_a\mu)(t)
=
a^{2k}\int_{\mathbb R}\varphi(au)\,d\mu(u)
\]
for every test function \(\varphi\).  In words, we scale the underlying variable by
\(u\mapsto au\), and multiply the mass by \(a^{2k}\).  This operation preserves
compact support and total mass zero, so \(\widetilde T_a\mu\in\mathcal M_0\).

For \(\mu,\nu\in\mathcal M_0\), the definition gives
\[
\begin{aligned}
\langle \widetilde T_a\mu,\widetilde T_a\nu\rangle_\beta
&=
-\frac12\iint_{\mathbb R^2}|u-v|^\beta\,
d(\widetilde T_a\mu)(u)\,d(\widetilde T_a\nu)(v)\\
&=
a^{4k}\left(
-\frac12\iint_{\mathbb R^2}|au-av|^\beta\,d\mu(u)\,d\nu(v)
\right)\\
&=
a^{4k+\beta}\langle\mu,\nu\rangle_\beta
=
a^p\langle\mu,\nu\rangle_\beta .
\end{aligned}
\]
In particular, if \(\mu\) belongs to the null space of the semidefinite form
\(\langle\cdot,\cdot\rangle_\beta\), then so does \(\widetilde T_a\mu\).  Hence
\(\widetilde T_a\) descends to the quotient space.  On the quotient, and then on
its Hilbert-space completion, we denote the induced operator by \(T_a\).  The
identity above gives
\[
\|T_a\xi\|_{\mathcal H_\beta}=a^{p/2}\|\xi\|_{\mathcal H_\beta}
\]
for every \(\xi\in\mathcal H_\beta\).

It remains to verify that this operator has the required action on the vectors
\(\Gamma_0(x)\).  By direct inspection of the definition of
\(\psi_x^{(0)}\), for every \(x\in[-1,1]\) and \(0<a\le1\),
\[
\psi_{ax}^{(0)}(t)=a^{2k-1}\psi_x^{(0)}(t/a).
\]
Equivalently, for every test function \(\varphi\),
\[
\int_{\mathbb R}\varphi(t)\psi_{ax}^{(0)}(t)\,dt
=
a^{2k}\int_{\mathbb R}\varphi(au)\psi_x^{(0)}(u)\,du.
\]
Moreover,
\(
\frac{|ax|^{2k}}{(2k)!}\delta_0
\)
is obtained from
\(
\frac{|x|^{2k}}{(2k)!}\delta_0
\)
by the same scaling rule, because the mass is multiplied by \(a^{2k}\) and the
point \(0\) remains fixed.  Therefore
\(
\nu_{ax}^{(0)}=\widetilde T_a\nu_x^{(0)}.
\)
Passing to equivalence classes in \(\mathcal H_\beta\), we obtain
\(
\Gamma_0(ax)=T_a\Gamma_0(x).
\)

Finally, we prove that \(\Gamma_0\) is continuous.  Applying
Proposition~\ref{prop:zero-anchored-exact-bridge-log}\textnormal{(i)} to the
two-point set \(\{x,y\}\) with \(\boldsymbol\lambda=(1,-1)\), we obtain
\[
|x-y|^p
=
c_{p,k}\|\Gamma_0(x)-\Gamma_0(y)\|_{\mathcal H_\beta}^2
+
\|M_{+,0}^{1/2}(\boldsymbol w_0(x)-\boldsymbol w_0(y))\|_2^2
-
\|M_{-,0}^{1/2}(\boldsymbol w_0(x)-\boldsymbol w_0(y))\|_2^2.
\]
Thus
\[
c_{p,k}\|\Gamma_0(x)-\Gamma_0(y)\|_{\mathcal H_\beta}^2
\le
|x-y|^p
+
\left|
\|M_{+,0}^{1/2}(\boldsymbol w_0(x)-\boldsymbol w_0(y))\|_2^2
-
\|M_{-,0}^{1/2}(\boldsymbol w_0(x)-\boldsymbol w_0(y))\|_2^2
\right|.
\]
Since \(\boldsymbol w_0\) is continuous on \([-1,1]\), the finite-dimensional correction tends
to \(0\) as \(y\to x\).  Hence
\[
\|\Gamma_0(x)-\Gamma_0(y)\|_{\mathcal H_\beta}\to0
\]
as \(y\to x\).  This completes the proof.



\end{proof}

\end{document}
